\section{Simultaneous selection of the auxiliary quadratic characters}
\label{qs-section}

The three quadratic characters used in the RH subcase do not describe
the entire requirement for Dirichlet GRH. For a fixed modulus one needs
a finite family of its characters, each twisted by either of two new
quadratic characters or by their product. The same two primes must work
throughout an interval of scales. This section proves that selection
from the following explicit analytic hypothesis:
\begin{equation}\label{qs-qrh}
 L(s,\chi)\ne0\qquad(\operatorname{Re}s>11/12)
\end{equation}
for every primitive Dirichlet character, with the pole of $\zeta$ at
$1$ allowed. This is the endpoint stated in the supplied October
quasi-Riemann-hypothesis paper; the September endpoint is stronger.
\cite{qrh-october,qrh-september}
The conclusion below remains conditional on \eqref{qs-qrh} in this
reconstruction. Its proof does not assume the Riemann hypothesis for
any of the selected characters.

For a nonprincipal Dirichlet character $\chi$, write
\[
 F_\chi(s)=-\frac{L'}{L}(s,\chi),\qquad
 G_\chi(y)=\sum_{n\le y}\Lambda(n)\chi(n)(y/n-1)-yF_\chi(1).
\]
Thus the second term in $G_\chi$ is
$+yL'(1,\chi)/L(1,\chi)$. For an odd prime $p$, let
$\chi_p(n)=(n/p)$ be the Legendre character, extended by zero at
multiples of $p$. In a product of these characters, a factor with
exponent zero is omitted; it is not a principal character at an
unnecessarily enlarged modulus. Each resulting primitive character
is considered at its own conductor.

\begin{theorem}[Simultaneous quadratic selection]\label{qs-selection}
Assume \eqref{qs-qrh}. Fix an integer $M\ge1$, real numbers
$0<\varepsilon<1/52$ and $A\ge1$. For all sufficiently large $Y$,
with threshold depending only on $M,\varepsilon,A$, put
\[
 z=Y^{\varepsilon/5},\qquad z'=z(\log Y)^2.
\]
There are distinct primes $q<d$, both coprime to $M$, such that
\begin{gather*}
 e^{z/2}<q<e^z,\quad q\equiv3\pmod8,\qquad
 e^{z'/2}<d<e^{z'},\quad d\equiv7\pmod{16},\\
 q\mid d+1,\qquad
 \ell\mid d+1\quad\hbox{for every odd prime }\ell\le z.
\end{gather*}
For every primitive character $\xi$ whose conductor divides $8M$,
and all $i,j\in\{0,1\}$ with $i+j>0$, the character
$\chi=\xi\chi_q^i\chi_d^j$ is primitive and nonprincipal, and
simultaneously for every $Y\le y\le Y^A$,
\begin{equation}\label{qs-bounds}
 |G_\chi(y)|\le y^{1/2+\varepsilon/2},\qquad
 \left|\frac{L'}{L}(1,\chi)\right|\le4\log y.
\end{equation}
Every such $L$-function is also zero-free in
$\operatorname{Re}s>399/400$, at every height.
\end{theorem}

The finite family of base characters is fixed before $Y$ varies.
The primes $q,d$, and consequently their quadratic twists, may change
with $Y$. This order of quantifiers is essential: the theorem gives
no square-root estimate for an arbitrary fixed character at every
large scale.

An imprimitive base character modulo $M$ is used through its primitive
inducer, whose conductor divides $M$. This supplies the primitive
character determined by its values away from $2Mqd$.
Primitive and imprimitive centered sums are not silently identified.
If a nonprincipal $\chi$ modulo $Q$ is induced from $\chi^*$ of
conductor $f$, direct subtraction of the finite Euler factors gives
the exact correction
\[
 G_\chi(y)-G_{\chi^*}(y)=
 \sum_{\substack{p\mid Q\\p\nmid f}}
   \sum_{r\ge1}(\log p)\chi^*(p)^r\min(1,y/p^r).
\]
Indeed, subtract the missing prime powers in the finite defining sum
and add their convergent geometric series in $-L'/L(1)$; the two
terms combine into the displayed minimum. The theorem is stated
at the primitive conductors, so these corrections are not needed
in its selection argument.

\subsection{Uniform analytic estimates with a smooth cutoff}

We record the conductor dependence used below. Under
\eqref{qs-qrh}, put $\sigma=23/24$. Uniformly for nonprincipal
characters modulo $Q$,
\begin{equation}\label{qs-logderivative}
 |F_\chi(v+it)|\ll\log\bigl((Q+3)(|t|+3)\bigr)
       \qquad(\sigma\le v\le2).
\end{equation}
Here and below the constants in analytic estimates are absolute once
the displayed lines and smooth functions have been fixed. For clarity,
this estimate does not mean that the value of $L$ is bounded below by
an absolute constant. The proof uses a holomorphic logarithm. For a
primitive nonprincipal character, its partial sums have modulus at
most $Q$, so partial summation gives
$|L(s,\chi)|\le Q|s|/\operatorname{Re}s$ in the right half-plane.
On a disk centered at $2+it$ with radius strictly between
$2-\sigma$ and $2-11/12$, take the logarithm whose value at the
center is supplied by the absolutely convergent Euler product.
That center value is bounded absolutely. The disk estimate of
Lemma~\ref{lem:qrh-log-disk}, followed by Cauchy's derivative
formula on a smaller disk, gives \eqref{qs-logderivative}.
The chosen radii are fixed, so no conductor enters its constant.

If $\chi$ is induced from $\chi^*$, its additional contribution is
\[
 \frac{L'}L(s,\chi)-\frac{L'}L(s,\chi^*)
 =\sum_{\substack{p\mid Q\\p\nmid\operatorname{cond}\chi^*}}
       \frac{\chi^*(p)\log p}{p^s-\chi^*(p)}.
\]
On the displayed strip this is at most
$(2^\sigma-1)^{-1}\log Q$ in modulus. The additional Euler
factors have their zeros on $\operatorname{Re}s=0$, so they do
not change the zero-free hypothesis in use. For a principal
character one instead bounds $F_\chi(s)-1/(s-1)$ by the same
quantity, using $(s-1)\zeta(s)$ in the disk argument.

Fix a smooth function $w:[0,\infty)\to[0,1]$ with
$w(u)=1$ for $u\le1$ and $w(u)=0$ for $u\ge2$.

\begin{lemma}[A finite smooth approximation]\label{qs-cutoff}
Under \eqref{qs-qrh}, for every nonprincipal character modulo $Q$
and $Z\ge2$,
\[
 F_\chi(1)=\sum_{n\le2Z}\frac{\Lambda(n)\chi(n)}n w(n/Z)
       +O\bigl(Z^{-1/24}\log(Q+3)\bigr).
\]
Consequently, uniformly for $1\le y\le Z$,
\begin{equation}\label{qs-smooth-identity}
 G_\chi(y)=-\sum_{n\le2Z}\Lambda(n)\chi(n)
                    \min(1,y/n)w(n/Z)
       +O\bigl(yZ^{-1/24}\log(Q+3)\bigr).
\end{equation}
There is also an absolute constant $C$ such that
\begin{equation}\label{qs-loglog}
 |F_\chi(1)|\le192\log\log(Q+3)+C.
\end{equation}
These statements exclude principal characters: $F_\chi(1)$ would
then have a pole.
\end{lemma}
\begin{proof}
For $\operatorname{Re}s>0$, the Mellin transform of $w$ is
\[
 \widehat w(s)=\int_0^\infty w(u)u^{s-1}\,du
       =-\frac1s\int_1^2 w'(u)u^s\,du.
\]
The second expression continues it meromorphically, with only a
simple pole at $0$, of residue $1$. Repeated integration by parts
after $u=e^v$ proves rapid decay on every fixed vertical strip away
from this pole. Mellin inversion and absolute convergence on
$\operatorname{Re}s=1$ give
\[
 \sum_n\frac{\Lambda(n)\chi(n)}n w(n/Z)
 =\frac1{2\pi i}\int_{(1)}F_\chi(1+s)\widehat w(s)Z^s\,ds.
\]
Move the contour to $\operatorname{Re}s=-1/24$.
Equation~\eqref{qs-logderivative} and rapid decay justify the shift
and bound the new integral by $O(Z^{-1/24}\log(Q+3))$.
The only crossed pole is that of $\widehat w$ at $0$, with residue
$F_\chi(1)$. This proves the first assertion with the stated sign.
When $n\le y\le Z$, $w(n/Z)=1$. Substitution in the definition
of $G_\chi$ therefore gives \eqref{qs-smooth-identity} exactly,
apart from the displayed analytic error.

We shall use the elementary bound
\begin{equation}\label{qs-chebyshev}
 \psi(t):=\sum_{n\le t}\Lambda(n)\le4t\qquad(t\ge1).
\end{equation}
Indeed, every prime power in $(m,2m]$ contributes its logarithm
to $\binom{2m}{m}$, and all other contributions to its prime
valuations are nonnegative. Thus
$\psi(2m)-\psi(m)\le2m\log2$.
Sum this for powers of two and use the smallest power of two
at least $t$ to obtain \eqref{qs-chebyshev}.
Partial summation now gives
\[
 \sum_{n\le T}\frac{\Lambda(n)}n\le4(1+\log T)
       \qquad(T\ge1).
\]
Apply the first assertion with $Z=(\log(Q+3))^{48}$.
The error is bounded absolutely, and the polynomial has norm at
most $4(1+\log(2Z))$. This is
$192\log\log(Q+3)+O(1)$, proving \eqref{qs-loglog}.
\end{proof}

\subsection{Quadratic orthogonality in a progression}

Choose a nonnegative $\Phi\in C_c^\infty(1/2,1)$ with
$I_\Phi=\int\Phi>0$. For $(a,m)=1$ and $U\ge2$, give a prime
$p\equiv a\pmod m$ weight $(\log p)\Phi(p/U)$. Let
\[
 W(U;m,a)=\sum_{p\equiv a\ (m)}(\log p)\Phi(p/U).
\]
Whenever $W>0$, expectation below means the weighted average with
these weights divided by $W$.

\begin{lemma}[Weighted correlations]\label{qs-correlations}
Assume \eqref{qs-qrh}. Uniformly in $m,a,U$,
\[
 W(U;m,a)=\frac{I_\Phi U}{\varphi(m)}
       +O\bigl(U^{23/24}\log(m+3)+U^{1/2}(\log(2U))^2\bigr).
\]
For odd squarefree $s>1$ coprime to $m$,
\[
 \left|\sum_{p\equiv a\ (m)}(\log p)\Phi(p/U)(p/s)\right|
 \ll U^{23/24}\log(ms+3)+U^{1/2}(\log(2U))^2.
\]
In particular, if $W\ge I_\Phi U/(2\varphi(m))$, then
for every such $s\le N^2$,
\begin{equation}\label{qs-rho}
 |\mathbb E(p/s)|\ll\varphi(m)
 \left(U^{-1/24}\log(mN^2+3)
             +U^{-1/2}(\log(2U))^2\right).
\end{equation}
\end{lemma}
\begin{proof}
For a Dirichlet character $\alpha$ modulo $L$, Mellin inversion gives
\[
 \sum_n\Lambda(n)\alpha(n)\Phi(n/U)
 =\frac1{2\pi i}\int_{(2)}F_\alpha(s)
                          \widehat\Phi(s)U^s\,ds.
\]
Moving the line to $23/24$ gives an error
$O(U^{23/24}\log(L+3))$ by \eqref{qs-logderivative}.
The pole at $1$ contributes $I_\Phi U$ precisely when $\alpha$
is principal. This argument includes imprimitive characters and
their local Euler factors, as described above.

Expand the progression condition with the characters modulo $m$.
For the first formula exactly one character is principal. For the
second, the products of these characters with $(\cdot/s)$ are all
nonprincipal modulo $ms$. To verify the latter assertion, the Chinese
remainder theorem permits an integer equal to $1$ modulo $m$ on
which the nontrivial character $(\cdot/s)$ has value $-1$.
The factor $1/\varphi(m)$ in the character average cancels the
number of error terms. Finally remove all prime powers of exponent
at least two. Their absolute total is
$O(U^{1/2}(\log(2U))^2)$, independently of the progression or
twist. Division by the lower bound for $W$ proves
\eqref{qs-rho}.
\end{proof}

The next lemma is elementary and independent of
\eqref{qs-qrh}. It makes explicit why forcing splitting at a short
initial set of primes does not destroy the second moment.

\begin{lemma}[A finite second moment]\label{qs-second-moment}
Let $\mathcal P$ be a finite probability space of primes $p>N$
in one reduced class $a\pmod m$, where $8\mid m$.
Let $2\le t\le N$, $1\le z\le t^{1/5}$, and suppose every prime
divisor of $m$ at most $N$ is at most $z$. Suppose that
\[
 |\mathbb E(p/s)|\le\rho
\]
for every odd squarefree $s\le N^2$, $s>1$, $(s,m)=1$.
Then, for every complex sequence $|\kappa(n)|\le1$,
\begin{equation}\label{qs-moment}
 \mathbb E\left|\sum_{n\le N}\Lambda(n)\min(1,t/n)
                         \kappa(n)(n/p)\right|^2
 \le C\left(t\log(2t)+\rho t^2\log^2(eN)\right)
\end{equation}
with an absolute constant $C$.
The same conclusion holds if the correlation hypothesis omits
at most one odd squarefree integer $s_0>1$.
\end{lemma}
\begin{proof}
Put $c_n=\Lambda(n)\min(1,t/n)$. Only prime powers occur.
For a fixed prime $\ell$, let
\[
 b_\ell=\sum_{j\ge1}c_{\ell^j},
\]
where the sum can be truncated at $\ell^j\le N$; all estimates
below remain valid for the infinite sum. If $\ell\le t$, a
geometric tail gives $b_\ell\le\log t+2\log\ell\le3\log t$.
If $\ell>t$, it gives
$b_\ell\le2t\log\ell/\ell$.
It follows that
\begin{equation}\label{qs-prime-diagonal}
 \sum_\ell b_\ell^2\ll t\log(2t),\qquad
 \max_\ell b_\ell\ll\log(2t).
\end{equation}
Here are the elementary details of the sum. From
\eqref{qs-chebyshev}, splitting the primes at $\sqrt t$ gives
$\pi(t)\le10t/\log t$ for $t\ge2$.
This bounds the terms $\ell\le t$. For the others, partial
summation with $\vartheta(u)=\sum_{\ell\le u}\log\ell\le4u$
gives
\[
 \sum_{\ell>t}\frac{(\log\ell)^2}{\ell^2}
 \le4\int_t^\infty\frac{2\log u-1}{u^2}\,du
 \ll\frac{\log(2t)}t.
\]
The claimed bounds follow, with a harmless adjustment on $2\le t<3$.

All even powers contribute at most
\[
 \sum_{r\ge1}\Lambda(r)\min(1,t/r^2)\le12\sqrt t
\]
by \eqref{qs-chebyshev} and partial summation at $\sqrt t$.
The total contribution from primes dividing $m$ is at most
$3z\log t$, since those contributing primes are at most $z$.
Because $z\le t^{1/5}$ and
$\log t\le(10/3)t^{3/10}$, these terms together have total
absolute value
\begin{equation}\label{qs-forced-part}
 R\ll\sqrt t.
\end{equation}

Quadratic reciprocity shows that $(\ell/p)$ is constant on the
given progression when $\ell\mid m$: the residue modulo $4$
is fixed, the residue modulo each odd such $\ell$ is fixed, and
the value at $2$ is fixed by the residue modulo $8$.
Group the sum in \eqref{qs-moment} into these terms and the even
powers, followed by odd powers of primes $\ell\nmid m$.
It becomes
\[
 A_0+\sum_{\ell\nmid m}A_\ell(\ell/p),
       \qquad |A_0|\le R,\quad |A_\ell|\le b_\ell.
\]
The diagonal quadratic factors equal $1$, since $p>N$.
The other factors are, up to a constant sign fixed by the
progression, $(p/\ell)$ or $(p/(\ell\ell'))$ for distinct
primes outside $m$. Expanding the square and using the correlation
bound therefore gives
\[
 \mathbb E|\cdots|^2\le R^2+\sum_\ell b_\ell^2
                 +\rho\left(\sum_{n\le N}c_n\right)^2.
\]
Partial summation gives
$\sum_{n\le N}c_n\le4t(1+\log(N/t))\le4t\log(eN)$.
Together with \eqref{qs-prime-diagonal} and
\eqref{qs-forced-part}, this proves the bound.

If there is one exceptional $s_0$, it affects these correlations
only when $s_0$ is one prime or the product of two distinct primes.
In the first case the extra term is at most
$2R\max b_\ell$; in the second it is at most
$2(\max b_\ell)^2$. Both are $O(t\log(2t))$.
Thus one exceptional correlation does not change the conclusion.
No multiplicativity of $\kappa$ has been used.
\end{proof}

\subsection{Uniformity on an interval of scales}

\begin{lemma}[A maximal second moment]\label{qs-maximal}
For the same finite probability space, put
\[
 S_p(t;\kappa)=\sum_{n\le N}\Lambda(n)\min(1,t/n)
                                      \kappa(n)(n/p).
\]
Suppose that, for every $Y\le t\le X$ and every
$|\kappa(n)|\le1$,
$\mathbb E|S_p(t;\kappa)|^2\le C t\log(2t)$, where
$2\le Y\le X\le N$. For each fixed such sequence and each
$\tau>0$, the probability that
\[
 |S_p(t;\kappa)|>\tfrac12t^{1/2+\tau}
       \quad\hbox{for some }t\in[Y,X]
\]
is at most
\begin{equation}\label{qs-bad-probability}
 16C Y^{-2\tau}
 \left\{\frac{\log(4Y)}{1-2^{-2\tau}}
       +\frac{(\log2)2^{-2\tau}}{(1-2^{-2\tau})^2}\right\}.
\end{equation}
For any fixed finite set of sequences, multiply this bound by
their number to control failure for at least one sequence.
\end{lemma}
\begin{proof}
The function $S_p(t;\kappa)$ is continuous and piecewise linear.
At all but finitely many points,
\[
 D_p(t):=S'_p(t;\kappa)
       =\sum_{t<n\le N}\frac{\Lambda(n)}n\kappa(n)(n/p).
\]
Its second moment is at most $C\log(2t)/t$: multiply by $t$
and apply the hypothesis to
$\kappa(n)1_{n>t}$. The uniformity in bounded sequences is what
permits this step.

On a block $[T,W]$ with $T\le W\le2T$, the fundamental theorem
of calculus and Cauchy--Schwarz give
\[
 \sup_{T\le t\le W}|S_p(t;\kappa)|^2
 \le2|S_p(T;\kappa)|^2+2T\int_T^W|D_p(v)|^2\,dv.
\]
Taking expectations yields a bound $4CT\log(4T)$.
Markov's inequality therefore bounds failure in that block by
$16CT^{-2\tau}\log(4T)$. Cover $[Y,X]$ by the blocks starting
at $T=2^jY$ and sum this bound. The two geometric sums are
$\sum r^j=(1-r)^{-1}$ and
$\sum jr^j=r(1-r)^{-2}$, with $r=2^{-2\tau}$.
This proves \eqref{qs-bad-probability}; the finite-family assertion
is the union bound. There is no discretization loss depending on
the length $X/Y$.
\end{proof}

\subsection{Proof of simultaneous selection}

\begin{proof}[Proof of Theorem~\ref{qs-selection}]
Let $\mathcal F$ be the finite set of primitive characters whose
conductors divide $8M$, and put
\[
 X=Y^A,\qquad Z=Y^{48(A+1)},\qquad N=\lfloor2Z\rfloor.
\]
Its cardinality depends on $M$ alone. We use the finite sums with
the bounded sequence $\kappa(n)=\xi(n)w(n/Z)$, and later also
with $\kappa(n)=\xi(n)\chi_q(n)w(n/Z)$.

First take $U_1=e^z$, $m=8$, $a=3$, and use the probability
measure preceding Lemma~\ref{qs-correlations}. Its total weight
is positive for all sufficiently large $Y$. Every prime in its
support is greater than $U_1/2$, hence greater than $N^2$ and
$8M$. For every fixed $L>0$, the right side of
\eqref{qs-rho} is at most $Y^{-L}$ for all sufficiently large
$Y$: $z=Y^{\varepsilon/5}$ grows faster than every multiple of
$\log Y$, whereas $N$ is a fixed power of $Y$.
This proves the correlation estimate uniformly for all the
integers needed in Lemma~\ref{qs-second-moment}.

For $Y\le t\le X$, the only prime divisor of $8$ is harmless
in \eqref{qs-forced-part}. More generally $z\le t^{1/5}$,
because $\varepsilon<1$. Choose $L=A+2$. The error term in
\eqref{qs-moment} is absorbed into $Ct\log(2t)$, uniformly
over this interval and over all bounded sequences. Apply
Lemma~\ref{qs-maximal} with $\tau=\varepsilon/2$ to all
$\xi\in\mathcal F$. Their total failure probability tends
to zero. We can consequently choose a prime $q$ in the support
such that all their finite sums have modulus at most
$\tfrac12t^{1/2+\varepsilon/2}$ throughout $[Y,X]$.
This prime satisfies $q\equiv3\pmod8$ and
$e^{z/2}<q<e^z$ once $Y$ is large.

Now put
\[
 P_z=\prod_{\substack{3\le\ell\le z\\\ell\ {
m prime}}}\ell,
       \qquad m=16qP_z,
\]
and choose by the Chinese remainder theorem the reduced residue
class $a$ given by
$a\equiv7\pmod{16}$ and $a\equiv-1\pmod{qP_z}$.
Take $U_2=e^{z'}$. The bound \eqref{qs-chebyshev} gives
\[
 \log m\le\log16+5z,
\]
uniformly in the previously chosen $q$.
The modulus satisfies the local condition in
Lemma~\ref{qs-second-moment}: the prime $q$ is greater than
$N$, and every other odd prime divisor is at most $z$.

Again the total prime weight is positive, and for every fixed
$L>0$ the correlation bound is at most $Y^{-L}$, uniformly in
$q$. Indeed its logarithm is bounded above, apart from
$O_{A,M,\varepsilon}(\log Y+\log z')$, by
\[
 5z-z'/24=5z-z(\log Y)^2/24,
\]
which is eventually smaller than $-L\log Y$. The prime-power
error has the still smaller exponent $5z-z'/2$.
All primes in the support exceed $N^2$ and $q$.

Apply the same second-moment and maximal estimates to the at
most $2|\mathcal F|$ sequences
$\xi(n)\chi_q(n)^i w(n/Z)$, $\xi\in\mathcal F$, $i=0,1$.
Their absolute values are at most one; the constants in the
second moment do not depend on the already chosen $q$.
The total failure probability again tends to zero. Choose $d$
outside this exceptional set. It has all the stated congruences
and sizes, and $q<d$. Both primes exceed $8M$.

The product of primitive characters of relatively prime
conductors is primitive: by the Chinese remainder theorem its
restriction to each conductor factor is the original primitive
character, so a character inducing the product must contain
each of those factors in its conductor. Applying this to
$\xi$, $\chi_q$, and $\chi_d$ shows that the conductor of each
selected twist is
$f q^i d^j$, where $f=\operatorname{cond}\xi\mid8M$.
It is nonprincipal, since at least one new prime occurs to the
first power in its conductor.

For all selected twists this conductor $Q$ satisfies
\[
 \log(Q+3)\ll_M z(\log Y)^2.
\]
Equation~\eqref{qs-smooth-identity} has error, uniformly on
$[Y,X]$, at most
\[
 C_M Y^A Y^{-2(A+1)}Y^{\varepsilon/5}(\log Y)^2=o(1).
\]
For large $Y$ this is at most
$\tfrac12Y^{1/2+\varepsilon/2}$.
The choices of $q,d$ therefore give the first bound in
\eqref{qs-bounds} for all twists simultaneously.

For the derivative bound, the same conductor estimate implies
\[
 \log\log(Q+3)
 \le\frac{\varepsilon}{5}\log Y+2\log\log Y+O_M(1).
\]
Equation~\eqref{qs-loglog} now yields
\[
 |F_\chi(1)|
 \le\frac{192\varepsilon}{5}\log Y
                 +384\log\log Y+O_M(1)
 \le4\log Y\le4\log y
\]
for sufficiently large $Y$. Here
$192\varepsilon/5<192/260<4$, leaving a fixed positive
margin. The zero-free assertion follows directly from
\eqref{qs-qrh}, since $399/400>11/12$.
\end{proof}

For the RH subcase take the principal base character of conductor
$1$. The three twists are precisely $\chi_q,\chi_d,\chi_q\chi_d$.
Since $q,d\equiv3\pmod4$, quadratic reciprocity identifies
$\chi_q(n)$ and $\chi_d(n)$ with the Kronecker characters
$(-q/n)$ and $(-d/n)$, including the supplementary rule at $2$.
For each odd prime $\ell\mid d+1$, $-d\equiv1\pmod\ell$,
so $\ell$ splits in $\mathbb Q(\sqrt{-d})$. The same holds for
$\ell=q$. At $2$, the congruence $d\equiv7\pmod{16}$ makes
$X^2-X+(d+1)/4$ have the two distinct roots $0,1$ modulo $2$.
Thus the required splitting statements are also supplied.

For ordinary RH, the selected character group has order one and the
translation subgroup is the whole group $A$. The subset-share condition
is then exact. For a nonnegative integer $e$, the number fields in this construction are
\[
 \mathbf Q(\zeta_{2^{e+2}})^+(\sqrt{-q},\sqrt{-d}),
 \qquad [E:\mathbf Q]=4\cdot2^e.
\]
Here $A$ is the Galois group of the real cyclotomic subfield, of order
$2^e$. The prescribed size requirement is
$x^{2+8\varepsilon}\le2^e\le2x^{2+8\varepsilon}$.
Thus the simultaneous selection above reduces the main auxiliary
family to three quadratics, but the field degrees and the ranks used in this construction still grow. Quantitative integral cohomology, integral theta
maps, period-normalized global matching and section-lattice degree
estimates remain to be supplied for that family.

The argument is uniform if $z'$ is allowed to vary between
$z(\log Y)^2$ and $A^2z(\log Y)^2$: its lower endpoint
controls correlations, and its upper endpoint controls
$\log(Q+3)$. No constant depends on the growing primes or on
the intermediate scale $y$.

\subsection{The larger-scale character estimate}

The same analytic input supplies the separate small relative error
needed on the much larger scales. For nonprincipal $\chi$ modulo
$Q$, Mellin inversion of $(u-1)_+$ gives
\[
 \sum_{n\le y}\Lambda(n)\chi(n)(y/n-1)
 =\frac1{2\pi i}\int_{(2)}
             F_\chi(s)\frac{y^s}{s(s-1)}\,ds.
\]
Moving the line to $23/24$ crosses the simple pole with residue
$yF_\chi(1)$. Subtracting it gives
\begin{equation}\label{qs-large-scale}
 |G_\chi(y)|\ll y^{23/24}\log(Q+3)\qquad(y\ge1).
\end{equation}
The vertical integral is absolutely convergent by
\eqref{qs-logderivative}. For the primitive principal character
one uses instead
\[
 G_1(y)=\sum_{n\le y}\Lambda(n)(y/n-1)
                   -y(\log y-1-\gamma).
\]
Since $-\zeta'/\zeta(s)=1/(s-1)-\gamma+O(s-1)$, the residue
of the double pole is exactly the term subtracted here, and the
same argument gives $|G_1(y)|\ll y^{23/24}$.
This treats the principal case separately rather than assigning
it a nonexistent logarithmic derivative at $1$.

For example, if $\log(Q+3)\ll x^{\varepsilon/5}(\log x)^2$
and $y\ge\tfrac13\exp(x^{\varepsilon/3})$, then
\eqref{qs-large-scale} gives, for all sufficiently large $x$,
\[
 |G_\chi(y)|\le y\exp(-x^{\varepsilon/8}).
\]
This follows because $y^{-1/24}$ decays faster than the displayed
polynomial conductor factor and than $\exp(-x^{\varepsilon/8})$.
Thus a fixed QRH gap is sufficient for this larger-scale input
as well as for the simultaneous selection, though it is
insufficient for a direct square-root estimate for every fixed
character.

\subsection{A necessary restriction and the remaining scope}

The restriction on prime divisors of the progression modulus in
Lemma~\ref{qs-second-moment} has real content. If it were removed,
take a large integer $t$, put
$m=16\prod_{\substack{3\le\ell\le t\\\ell\ {\rm prime}}}\ell$,
and choose the reduced class
$a\equiv7\pmod{16}$, $a\equiv-1\pmod\ell$ for every such
odd prime. Quadratic reciprocity and the supplementary rule at
$2$ give $(n/p)=1$ for every $n\le t$ and every prime $p>t$
in this class. With $N=t$ and $\kappa=1$, the sum is therefore
$\psi(t)$ for every $p$. This is $\gg t$: indeed
\[
 \psi(2r)\ge\log\binom{2r}{r}
       \ge2r\log2-\log(2r+1),
\]
and monotonicity handles intervening real $t$.
The second moment is consequently $\gg t^2$, contradicting
$O(t\log t)$. Under \eqref{qs-qrh} the progression contains
primes for all sufficiently large averaging scales, so the
example is nonempty. This is a counterexample to dropping a
hypothesis, not to the actual selected-prime construction.

The finite second-moment and maximal lemmas have been proved here
without QRH. QRH is used to establish the weighted prime
correlations uniformly in the growing progression modulus, to
control the smooth logarithmic-derivative remainder, and to make
all required zero-free boxes automatic. The allowance of one
exceptional correlation also explains why the ordinary
exceptional-character formulation can be handled in the original
second moment; proving its full analytic input without QRH is a
separate route.

This section completes the selected-character part conditional on
the stated QRH theorem. It does not construct or estimate the
automorphic interpolation maps, their local minors, or their
nearly isotropic complex realizations. Nor does it supply primes
in varying quadratic ray classes. Those parts of the proposed
GRH proof remain separate obligations.


